\section{Motivation}

\begin{frame}{Warum 6-Achs-Druck?}
    \begin{columns}[T]
        \column{0.6\textwidth}
            \textbf{Limitierungen klassischer 3-Achs-Verfahren:}
            \begin{itemize}
                \item Schichten müssen parallel zur Bauplattform sein.
                \item Treppenstufeneffekte auf gekrümmten Oberflächen.
                \item Hoher Bedarf an Stützstrukturen bei Überhängen.
                \item Mechanische Schwäche entlang der Schichtgrenzen.
            \end{itemize}
            
            \vspace{0.25cm}
            \textbf{Potential von 6 Achsen:}
            \begin{itemize}
                \item \textcolor{SecondaryColor}{\textbf{Nichtplanares Slicen:}} 
                    Erhöhte Stabilität durch nicht-parallele und nicht-planare Schichten.
                \item \textcolor{SecondaryColor}{\textbf{Support-Reduktion:}} 
                    Überhänge ohne Stütze durch Änderung der Druckrichtung.
                \item \textcolor{SecondaryColor}{\textbf{Oberflächenqualität:}}
                    Glattere Oberflächen durch Anpassung der Schichtgeometrie an die Modellkonturen.
            \end{itemize}

        \column{0.4\textwidth}
            % Vergleichsgrafik: Treppe vs Smooth
            \centering
            \begin{tcolorbox}[
                enhanced, colback=white, colframe=PrimaryColor!30, arc=2mm, boxrule=0.5pt,
                title={Planar (3-Achs)}, fonttitle=\bfseries\tiny, coltitle=PrimaryColor, colbacktitle=PrimaryColor!10
            ]
                \centering
                \begin{tikzpicture}[scale=0.5]
                    \draw[gray!50] (0,0) -- (3,1.5); % Slope
                    \foreach \x in {0,0.5,...,2.5} {
                        \draw[fill=PrimaryColor!40] (\x, {0.5*\x}) rectangle (\x+0.6, {0.5*\x+0.25});
                    }
                \end{tikzpicture}
            \end{tcolorbox}
            
            \vspace{0.2cm}

            \begin{tcolorbox}[
                enhanced, colback=white, colframe=SecondaryColor!30, arc=2mm, boxrule=0.5pt,
                title={Nichtplanar (6-Achs)}, fonttitle=\bfseries\tiny, coltitle=SecondaryColor, colbacktitle=SecondaryColor!10
            ]
                \centering
                \begin{tikzpicture}[scale=0.5]
                    \draw[fill=SecondaryColor!40, draw=SecondaryColor] (0,0) -- (3,1.5) -- (3,2.0) -- (0,0.5) -- cycle;
                \end{tikzpicture}
            \end{tcolorbox}
    \end{columns}
\end{frame}
